\documentclass[letterpaper]{article} 
\usepackage[utf8]{inputenc}
\linespread{0.85}
\usepackage[T1]{fontenc}
\usepackage{amsmath}
\usepackage{amsfonts}
\usepackage{amssymb}
\usepackage{array}
\usepackage{fontspec}
\usepackage{booktabs}
\usepackage{hyperref}
\usepackage[version=4]{mhchem}
\usepackage{stmaryrd}
\usepackage[dvipsnames]{xcolor}
\colorlet{LightRubineRed}{RubineRed!70}
\colorlet{Mycolor1}{green!10!orange}
\definecolor{Mycolor2}{HTML}{00F9DE}
\usepackage{graphicx}
\usepackage{amsmath}
\usepackage{graphicx}
\usepackage{capt-of}
\usepackage{lipsum}
\usepackage{fancyvrb}
\usepackage{tabularx}
\usepackage{listings}
\usepackage[export]{adjustbox}
\graphicspath{ {./images/} }
\usepackage[utf8]{inputenc}
\usepackage[english]{babel}
\usepackage{float}
\usepackage{lipsum}
\usepackage{graphicx}
\usepackage{float}
\usepackage[margin=0.7in]{geometry}
\usepackage{amsmath}
\usepackage{graphicx}
\usepackage{capt-of}
\usepackage{tcolorbox}
\usepackage{lipsum}
\usepackage{graphicx}
\usepackage{float}
\usepackage{listings}
\definecolor{deepred}{RGB}{139,0,0}
\usepackage{hyperref} 
\usepackage{xcolor} % For custom colors
\lstset{
	language=Python,                % Choose the language (e.g., Python, C, R)
	basicstyle=\ttfamily\small, % Font size and type
	keywordstyle=\color{blue},  % Keywords color
	commentstyle=\color{gray},  % Comments color
	stringstyle=\color{red},    % String color
	numbers=left,               % Line numbers
	numberstyle=\tiny\color{gray}, % Line number style
	stepnumber=1,               % Numbering step
	breaklines=true,            % Auto line break
	backgroundcolor=\color{black!5}, % Light gray background
	frame=single,               % Frame around the code
}
\usepackage{float}
\usepackage[]{amsthm} %lets us use \begin{proof}
	\usepackage[]{amssymb} %gives us the character \varnothing
	
	\title{Assignment 3, MATH 4155}
	\author{Zongyi Liu}
	\date{Wed, Feb 25, 2026}
	\begin{document}
		\maketitle
		
		\section{Question 1}
		
		{\color{deepred}\fontspec{Georgia} HW6 Q1 Lyapunov Inequality}
		
		If $X$ is a measurable function on a probability space $(\Omega, \mathcal{F}, \mathbb{P})$ and $0<p<q<\infty$, then
		
		$$
		\|X\|_{p} \leq\|X\|_{q}
		$$
		
		In particular, for $0<p<q<\infty$ we have $\mathbb{L}^{q} \subseteq \mathbb{L}^{p}$ on a probability space. (This kind of inclusion can fail, and rather dramatically, on a general measure space.)
		
		
		
		\textbf{Answer}
		
			Let $0<p<q<\infty$ and suppose $\|X\|_q<\infty$. Then
			$\|X\|_p^p=\mathbb E(|X|^p)$.
			Set $r=\frac{q}{p}>1$ and $r'=\frac{q}{q-p}$ so that $\frac1r+\frac1{r'}=1$.
			Applying Hölder's inequality to $|X|^p$ and $1$ gives
			$
			\mathbb E(|X|^p)
			=\mathbb E(|X|^p\cdot 1)
			\le
			\left(\mathbb E(|X|^{pr})\right)^{1/r}
			\left(\mathbb E(1^{r'})\right)^{1/r'} .
			$
			
			Since $pr=q$ and $\mathbb P(\Omega)=1$, we obtain
			$
			\mathbb E(|X|^p)
			\le
			\left(\mathbb E(|X|^q)\right)^{1/r}
			=
			\left(\mathbb E(|X|^q)\right)^{p/q}.
			$ Taking the $p$-th root yields
			$
			\|X\|_p
			=
			\left(\mathbb E(|X|^p)\right)^{1/p}
			\le
			\left(\mathbb E(|X|^q)\right)^{1/q}
			=
			\|X\|_q .
			$ 
			Hence $\|X\|_p\le \|X\|_q$. In particular, if $X\in\mathbb L^q$ then
			$\|X\|_q<\infty$ implies $\|X\|_p<\infty$, so $\mathbb L^q\subseteq\mathbb L^p$
			on a probability space.
		
		\clearpage
		
		\section{Question 2}
		
		{\color{deepred}\fontspec{Georgia} HW6 Q3 An Extended Monotone Convergence Theorem}
		
		Suppose $X_{1} \leq \cdots \leq X_{n} \leq X_{n+1} \leq \cdots$ are integrable random variables, monotonically increasing pointwise to the random variable $X=\lim _{n \rightarrow \infty} X_{n}$.
		
		If $\sup _{n \in \mathbb{N}} \mathbb{E}\left(X_{n}\right)<\infty$, show that $X$ is integrable and that we have
		
		$$
		\lim _{n \rightarrow \infty} \mathbb{E}\left(X_{n}\right)=\mathbb{E}(X)
		$$
		
		\textbf{Answer}
		
		Since $X_1 \le X_n \le X$ for all $n$, we have
		$0 \le X_n - X_1 \uparrow X - X_1$ pointwise. Let $Y_n := X_n - X_1$. Then $Y_n \ge 0$ and $Y_n \uparrow X - X_1$.  
		By the Monotone Convergence Theorem,
		$\mathbb{E}[X - X_1] = \mathbb{E}[\lim_{n\to\infty} (X_n - X_1)]
		= \lim_{n\to\infty} \mathbb{E}[X_n - X_1]
		= \lim_{n\to\infty}(\mathbb{E}[X_n] - \mathbb{E}[X_1])$.
		
		Because $\sup_{n\in\mathbb{N}} \mathbb{E}[X_n] < \infty$, the right-hand side is finite, hence
		$\mathbb{E}[X - X_1] < \infty$. Since $X_1$ is integrable and $|X| \le |X_1| + (X - X_1)$, we obtain
		$\mathbb{E}[|X|] \le \mathbb{E}[|X_1|] + \mathbb{E}[X - X_1] < \infty$.
		Thus $X \in L^1$. Finally, again by the Monotone Convergence Theorem (MCT) applied to $X_n - X_1 \uparrow X - X_1$,
		$\lim_{n\to\infty} \mathbb{E}[X_n - X_1] = \mathbb{E}[X - X_1]$. Therefore
		$\lim_{n\to\infty} \mathbb{E}[X_n]
		= \mathbb{E}[X - X_1] + \mathbb{E}[X_1]
		= \mathbb{E}[X]$.
			
		
		\clearpage
		
		\section{Question 3}
		
		{\color{deepred}\fontspec{Georgia} HW6 Q4 Riemann-Zeta Calculus} 
		
		For $s>1$ fixed, let $X:\Omega\to\mathbb{N}$ be a random variable with the ``zeta distribution''
		\[
		\mathbb{P}(X=n)=\frac{1}{\zeta(s)\cdot n^{s}}, \qquad n\in\mathbb{N}, 
		\qquad \text{where} \qquad 
		\zeta(s):=\sum_{n\in\mathbb{N}}\frac{1}{n^{s}}
		\]
		is the Riemann zeta-function. For each $m\in\mathbb{N}$, consider the event
		\[
		A_m:=\{\omega\in\Omega: m \text{ is a factor of } X(\omega)\}
		=\{\omega\in\Omega: X(\omega)=m k(\omega), \text{ for some } k(\omega)\in\mathbb{N}\}.
		\]
		
		\begin{enumerate}
			\item[(i)] Show that $\mathbb{P}(A_m)=1/m^{s}$.
			
			\item[(ii)] With $\mathfrak{P}$ denoting the set of prime numbers, conclude from (i) that the events in the collection $\{A_p\}_{p\in\mathfrak{P}}$ are independent,
			
			\item[(iii)] Use (i) and (ii) to derive the celebrated Euler formula
			\[
			\frac{1}{\zeta(s)}=\prod_{p\in\mathfrak{P}}\left(1-\frac{1}{p^{s}}\right).
			\]
		\end{enumerate}
		
		\textbf{Answer}
		
		\begin{enumerate}
			
			\item[(i)] We compute directly, and can get the result that
			$
			\mathbb{P}(A_m)
			=
			\sum_{k=1}^{\infty}\mathbb{P}(X=mk)
			=
			\sum_{k=1}^{\infty}\frac{1}{\zeta(s)(mk)^s}
			=
			\frac{1}{\zeta(s)m^s}\sum_{k=1}^{\infty}\frac{1}{k^s}
			=
			\frac{1}{\zeta(s)m^s}\zeta(s)
			=
			\frac{1}{m^s}.
			$
			
			\item[(ii)] Let $p_1,\dots,p_r$ be distinct primes. Then
			$
			A_{p_1}\cap\cdots\cap A_{p_r}
			=
			A_{p_1\cdots p_r},
			$
			since a positive integer is divisible by each $p_j$ iff it is divisible by their product.
			Thus by (i),
			$
			\mathbb{P}(A_{p_1}\cap\cdots\cap A_{p_r})
			=
			\mathbb{P}(A_{p_1\cdots p_r})
			=
			\frac{1}{(p_1\cdots p_r)^s}
			=
			\prod_{j=1}^r \frac{1}{p_j^s}
			=
			\prod_{j=1}^r \mathbb{P}(A_{p_j}).
			$
			Hence $\{A_p\}_{p\in\mathfrak P}$ are independent.
			
			\item[(iii)] For any finite set of primes $F$,
			$
			\mathbb{P}\!\left(\bigcap_{p\in F}A_p^c\right)
			=
			\prod_{p\in F}\mathbb{P}(A_p^c)
			=
			\prod_{p\in F}\left(1-\frac{1}{p^s}\right),
			$
			with (i) and independence. Then we let $F_N=\{p\in\mathfrak P: p\le N\}$. Then the events decrease and
			$
			\bigcap_{p\in\mathfrak P}A_p^c
			=
			\{X\text{ divisible by no prime}\}
			=
			\{X=1\}.
			$
			
			By continuity from above,
			$
			\mathbb{P}(X=1)
			=
			\lim_{N\to\infty}
			\mathbb{P}\!\left(\bigcap_{p\in F_N}A_p^c\right)
			=
			\lim_{N\to\infty}
			\prod_{p\le N}\left(1-\frac{1}{p^s}\right).
			$
			But
			$
			\mathbb{P}(X=1)=\frac{1}{\zeta(s)}.
			$
			Therefore
			$
			\frac{1}{\zeta(s)}
			=
			\prod_{p\in\mathfrak P}\left(1-\frac{1}{p^s}\right).
			$
			
		\end{enumerate}
		
		\clearpage
		
		\section{Question 4}
		
		{\color{deepred}\fontspec{Georgia} HW6 Q9 Modes of Convergence of Random Variables}
		
		On a suitable probability space, construct
		\begin{enumerate}
			\item[(i)] a sequence of random variables that converges in probability, but not almost surely;
			
			\item[(ii)] a sequence of random variables that converges almost surely, but does not exhibit fast convergence in probability; and
			
			\item[(iii)] a sequence of random variables that converges almost surely, but not in $\mathbb{L}^{1}$.
		\end{enumerate}
		
		(\emph{Hint}: All these sequences can be constructed on the space $\Omega=(0,1]$ with Lebesgue measure $\mathbb{P}$ on its Borel sets)
		
		
		\textbf{Answer} 
		
		\begin{enumerate}
			
			\item[(i)] \emph{Convergence in probability, but not almost surely.}
			
			Let $\Omega=(0,1]$ with $\mathbb{P}$ equal to Lebesgue measure, and define
			$
			X_n(\omega):=\mathbf{1}_{I_n}(\omega),
			$
			where the intervals $I_n$ are arranged as follows. For each $m\geq 0$ and $1\leq k\leq 2^m$, let
			$
			I_{2^m+k-1}:=\left(\frac{k-1}{2^m},\frac{k}{2^m}\right].
			$
			This is the usual “typewriter sequence.”
			
			For each $n$, the interval $I_n$ has length
			$
			\mathbb{P}(I_n)=2^{-m}
			$
			for some $m\to\infty$, hence
			$
			\mathbb{P}(|X_n|>\varepsilon)=\mathbb{P}(X_n=1)=\mathbb{P}(I_n)\to 0
			$
			for every $\varepsilon\in(0,1]$. Thus $X_n\to 0$ in probability.
			
			However, for each fixed $\omega\in(0,1]$ and each $m$, there is exactly one dyadic interval of length $2^{-m}$ containing $\omega$, so along the block
			$
			n=2^m,\dots,2^{m+1}-1
			$
			there is exactly one index with $X_n(\omega)=1$. Hence for every $\omega$,
			$
			X_n(\omega)=1 \text{ infinitely often, and } X_n(\omega)=0 \text{ infinitely often}.
			$
			So $X_n(\omega)$ does not converge for any $\omega\in(0,1]$. Therefore $X_n$ does {not} converge almost surely.
			
			\item[(ii)] \emph{Almost sure convergence, but not fast convergence in probability.}
			
			Define
			$
			X_n(\omega):=\mathbf{1}_{(0,1/n]}(\omega).
			$
			Then for every $\omega>0$, we have $X_n(\omega)=0$ for all sufficiently large $n$, because eventually $1/n<\omega$. Hence
			$
			X_n(\omega)\to 0
			$
			for every $\omega\in(0,1]$, so $X_n\to 0$ almost surely. On the other hand, for $\varepsilon\in(0,1]$,
			$
			\mathbb{P}(|X_n|>\varepsilon)=\mathbb{P}(X_n=1)=\frac{1}{n}.
			$
			Thus
			$
			\sum_{n=1}^\infty \mathbb{P}(|X_n|>\varepsilon)
			=
			\sum_{n=1}^\infty \frac{1}{n}
			=
			\infty.
			$
			So the convergence is not fast in probability.
			
			\item[(iii)] \emph{Almost sure convergence, but not convergence in $\mathbb{L}^1$.}
			
			Define
			$
			X_n(\omega):=n\,\mathbf{1}_{(0,1/n]}(\omega).
			$
			Again, for each fixed $\omega>0$, we have $X_n(\omega)=0$ for all sufficiently large $n$, so
			$
			X_n(\omega)\to 0
			$
			for every $\omega\in(0,1]$. Hence $X_n\to 0$ almost surely. But
			$
			\mathbb{E}[|X_n|]
			=
			\int_0^1 n\,\mathbf{1}_{(0,1/n]}(\omega)\,d\omega
			=
			n\cdot \frac{1}{n}
			=
			1
			$
			for all $n$. Therefore
			$
			\mathbb{E}[|X_n-0|]=1 \not\to 0,
			$
			so $X_n$ does not converge to $0$ in $\mathbb{L}^1$.
			
		\end{enumerate}
		
		
		\clearpage
		
		\section{Question 5}
		
		{\color{deepred}\fontspec{Georgia} HW6 Q12 The Concentration of Measure Phenomenon}
			
			\begin{itemize}
				
				\item[(i)] Suppose $X_1,X_2,\cdots$ are square-integrable random variables with
				$m=\mathbb{E}(X_j)$,
				$\mathbb{E}[(X_i-m)(X_j-m)]=0$ for $i\ne j$, and
				$\operatorname{Var}(X_j)=\mathbb{E}[(X_j-m)^2]\le C$
				for all $j$ and some constants $m$ and $C>0$.
				
				Show that, both in $\mathbb{L}^2$ and in probability,
				\[
				\lim_{n\to\infty}\frac{1}{n}\sum_{j=1}^n X_j = m.
				\]
				
				\item[(ii)] Suppose now that $Y_1,Y_2,\cdots$ are independent random variables with
				common distribution uniform on $(-1,1)$.
				Argue that
				\[
				\lim_{n\to\infty}\frac{1}{n}\sum_{j=1}^n Y_j^2=\frac{1}{3}
				\]
				holds in probability. (Does this hold also in a stronger sense?)
				
				\item[(iii)] Consider the cube $\mathbb{K}^n := (-1,1)^n$ in $\mathbb{R}^n$.
				Justify the statement:
				
				\emph{``For $n$ large, most of the volume of this cube is concentrated very
					near the sphere of radius $\sqrt{n/3}$ in $\mathbb{R}^n$.''}
				
				Establish for every $\varepsilon\in(0,1)$ that
				\[
				\lim_{n\to\infty}
				\frac{1}{2^n}\lambda_n\!\left(A_{n,\varepsilon}\cap\mathbb{K}^n\right)=1,
				\]
				where $\lambda_n$ denotes Lebesgue measure in $\mathbb{R}^n$ and
				$\lambda_n(\mathbb{K}^n)=2^n$, and
				\[
				A_{n,\varepsilon}
				=
				\left\{
				y\in\mathbb{R}^n:
				\frac{n}{3}(1-\varepsilon)
				\le
				y_1^2+\cdots+y_n^2
				\le
				\frac{n}{3}(1+\varepsilon)
				\right\}.
				\]
				
			\end{itemize}
		
		is a ``thin crust'' around the sphere of radius $\sqrt{n/3}$.
		
		\textbf{Answer} 
		
		\begin{enumerate}
			
			\item[(i)] Let
			$
			\overline{X}_n:=\frac{1}{n}\sum_{j=1}^n X_j.
			$
			We show that $\overline{X}_n\to m$ in $\mathbb{L}^2$, hence also in probability. First,
			$
			\mathbb{E}(\overline{X}_n)=\frac{1}{n}\sum_{j=1}^n \mathbb{E}(X_j)=m.
			$
			So it suffices to estimate the variance:
			$
			\mathbb{E}\bigl[(\overline{X}_n-m)^2\bigr]
			=
			\operatorname{Var}(\overline{X}_n).
			$
			Now
			$
			\overline{X}_n-m=\frac{1}{n}\sum_{j=1}^n (X_j-m),
			$
			hence
			$
			\operatorname{Var}(\overline{X}_n)
			=
			\mathbb{E}\!\left[\left(\frac{1}{n}\sum_{j=1}^n (X_j-m)\right)^2\right]
			=
			\frac{1}{n^2}\sum_{i=1}^n\sum_{j=1}^n
			\mathbb{E}\bigl[(X_i-m)(X_j-m)\bigr].
			$
			By assumption, the mixed terms vanish for $i\ne j$, so
			$
			\operatorname{Var}(\overline{X}_n)
			=
			\frac{1}{n^2}\sum_{j=1}^n \operatorname{Var}(X_j)
			\le
			\frac{1}{n^2}\sum_{j=1}^n C
			=
			\frac{C}{n}.
			$
			Therefore
			$
			\mathbb{E}\bigl[(\overline{X}_n-m)^2\bigr]\le \frac{C}{n}\to 0,
			$
			which proves $\overline{X}_n\to m$ in $\mathbb{L}^2$.
			Since $\mathbb{L}^2$ convergence implies convergence in probability
			(by Chebyshev's inequality: for every $\varepsilon>0$,
			$
			\mathbb{P}(|\overline{X}_n-m|>\varepsilon)
			\le
			\frac{\mathbb{E}[|\overline{X}_n-m|^2]}{\varepsilon^2}
			=
			\frac{\operatorname{Var}(\overline{X}_n)}{\varepsilon^2}
			\le
			\frac{C}{n\varepsilon^2}
			\to 0),
			$
			we also conclude $\overline{X}_n\to m$ in probability.
			
			\item[(ii)] Let $X_j:=Y_j^2$. Since $Y_j\sim \mathrm{Unif}(-1,1)$, the variables $X_j$ are independent and identically distributed, and
			$
			\mathbb{E}(X_j)=\mathbb{E}(Y_j^2)
			=
			\frac{1}{2}\int_{-1}^1 y^2\,dy
			=
			\frac{1}{2}\left[\frac{y^3}{3}\right]_{-1}^{1}
			=
			\frac{1}{2}\cdot\frac{2}{3}
			=
			\frac{1}{3}.
			$
			Also $0\le X_j\le 1$, so in particular $X_j\in\mathbb{L}^2$.
			
			Thus part (i) applies with $m=\frac{1}{3}$, and gives
			$
			\frac{1}{n}\sum_{j=1}^n Y_j^2
			=
			\frac{1}{n}\sum_{j=1}^n X_j
			\to
			\frac{1}{3}
			$
			in $\mathbb{L}^2$, hence in probability. In fact, this holds in a stronger sense: since $(X_j)$ are i.i.d.\ and integrable, the Strong Law of Large Numbers yields
			$
			\frac{1}{n}\sum_{j=1}^n Y_j^2 \to \frac{1}{3}
			$
			almost surely.
			
			\item[(iii)] Let $Y_1,\dots,Y_n$ be i.i.d. uniform on $(-1,1)$. Then the random vector
			$
			Y=(Y_1,\dots,Y_n)
			$
			is uniformly distributed on $\mathbb{K}^n=(-1,1)^n$, with density
			$
			\frac{1}{2^n}\mathbf{1}_{\mathbb{K}^n}(y).
			$
			Hence for any Borel set $B\subseteq \mathbb{R}^n$,
			$
			\mathbb{P}(Y\in B)=\frac{1}{2^n}\lambda_n(B\cap \mathbb{K}^n).
			$
			
			Now
			$
			Y\in A_{n,\varepsilon}
			$
			iff
			$
			\frac{n}{3}(1-\varepsilon)
			\le
			Y_1^2+\cdots+Y_n^2
			\le
			\frac{n}{3}(1+\varepsilon),
			$
			that is,
			$
			Y\in A_{n,\varepsilon}
			$
			iff
			$
			\left|
			\frac{1}{n}\sum_{j=1}^n Y_j^2-\frac13
			\right|
			\le
			\frac{\varepsilon}{3}.
			$
			Therefore
			$
			\frac{1}{2^n}\lambda_n(A_{n,\varepsilon}\cap \mathbb{K}^n)
			=
			\mathbb{P}(Y\in A_{n,\varepsilon})
			=
			\mathbb{P}\!\left(
			\left|
			\frac{1}{n}\sum_{j=1}^n Y_j^2-\frac13
			\right|
			\le
			\frac{\varepsilon}{3}
			\right).
			$
			By part (ii), $\frac1n\sum_{j=1}^n Y_j^2\to \frac13$ in probability, so the right-hand side tends to $1$. Hence
			$
			\lim_{n\to\infty}
			\frac{1}{2^n}\lambda_n(A_{n,\varepsilon}\cap \mathbb{K}^n)=1.
			$
			
			This proves the quoted statement: for large $n$, with probability close to $1$ under the uniform distribution on $\mathbb{K}^n$, one has
			$
			Y_1^2+\cdots+Y_n^2\approx \frac{n}{3}.
			$
			So most of the volume of the cube lies in the thin spherical shell where the radius satisfies
			$
			\|y\|\approx \sqrt{\frac{n}{3}}.
			$
			
		\end{enumerate}
		
		
	\end{document}
